\chapter{実験結果}

\section{学習結果}

\figfull[.96\textwidth]{baseline_training_loss.pdf}{本番運用モデルの全学習履歴。薄線は各記録、濃線は移動中央値、橙破線は現場運用のステップ19,000である。誤差低下は学習データへの適合を示すが、実機成功率ではない。}

学習損失は0.067703から0.000671へ低下し、最小記録は0.000499であった。初期に急減し、後半も緩やかに低下しており、学習データに対する目的関数値が下がり、教示動作への適合が進んだことを示す。未見条件での精度や実機成功率の向上は、この曲線だけでは判断できない。\ev{E08} 運用保存点は履歴の中盤にあり、損失最小点ではない。

\plain{この曲線が示すのは「人の手本に近い動作を出す学習が進んだ」ことまでである。「新しい机上配置でも成功する」ことを示すには、独立試験と1回ごとの成否記録が必要である。}

\section{データ規模と品質}

本番学習用25件の重複しないデータ保管単位はすべて読み取り可能で、1,051軌跡と879,852フレームの情報は整合した。3台の学習カメラ、14次元状態、14次元動作の項目構成も全データ保管単位で一致した。全数値検査では非有限値、全要素がゼロの動作レコード、軌跡内時刻不連続、フレーム番号欠落、完全重複軌跡は検出されなかった。\ev{E04,E05}

これは数値パイプラインの整合性を支持するが、全映像の完全性、意味ラベルの正しさ、データ漏洩がないことまで証明しない。試験分割がないため、学習・試験間の重複は現在評価できない。

\section{条件網羅と改善対象}

\figfull[.94\textwidth]{prompt_condition_gap.pdf}{タスク指示条件の網羅。キャップなし6件のデータ保管単位、158軌跡のいずれにも「キャップなしなら開栓を省略」という明示がなく、空回し改善に向けた具体的なデータ課題が分かる。}

キャップなし軌跡は存在するが、該当データ保管単位でも無条件の開栓指示が使われている。「キャップがある場合のみ」と明示したのは1件のデータ保管単位、19軌跡だけである。\ev{E11} これは、キャップなし画像と開栓継続の教師動作を同時に学ぶ可能性を示し、現場の空回し現象と整合する。

\limitbox{これは明確なデータ・指示条件のリスクだが、完全な因果証明ではない。同一ボトル・同一位置でキャップ有無を対にして試験し、開栓工程へ移行したかを記録する必要がある。}

方向関連404軌跡、反転条件127軌跡があり、向きに対する収集作業は既に行われている。一方、フレーム単位の向きラベルと条件別自律試験がないため、逆向き把持が視覚判断、把持対象選択、動作実行のどこで起きるかは未分離であり、発生率も算出できない。

\section{現場結果}

\begin{table}[H]
\centering
\caption{現場結果とエビデンス区分}
\begin{tabularx}{\textwidth}{>{\bfseries}p{42mm}p{34mm}X}
\toprule
観測内容 & 現段階の表現 & 表現を限定する理由\\
\midrule
把持、開栓、分別投入の全工程 & 現場展示で観測 & 観測記録と運用経路はあるが完全映像なし\\
連続稼働時間 & 1時間超（概算） & 開始終了時刻と中断台帳なし\\
平均処理速度 & 約2本/分（概算） & 逐次計数と有効稼働時間の定義なし\\
初期的な一般化 & ランダム位置と一定条件変化への適応兆候 & 固定条件、反復回数、未見ボトル一覧なし\\
安定成功率 & 未記録 & 学習損失や概算処理量から逆算しない\\
\bottomrule
\end{tabularx}
\end{table}

以上は、単発動作ではなく全工程と長時間反復が現場で観測されたことと整合する。しかし、一定の成功率や全ボトル型への一般化を支持するものではない。この境界が本報告の最も重要な解釈条件である。

\section{アテンション記録}

8,223件のアテンション記録における視点別構成比の中央値は、上部23.2\%、左手首33.2\%、右手首43.5\%であった。ボトル本体、口部、グリッパへ接近する工程で手首視点に相対的に大きな重みが割り当てられる傾向は、精密操作の必要性と整合する。\ev{E12} ただし成否ラベルがないため、色の強さからキャップ有無の理解や失敗原因を断定しない。

\section{強化学習のオフライン結果}

\figfull[.91\textwidth]{rlt_offline_validation.pdf}{強化学習連続ラウンドのオフライン検証。動作誤差は総じて低下したが、価値推定誤差は変動している。学習閉ループの成立は示すが、実機での改善は同条件比較を要する。}

ラウンド28の初期記録からラウンド32の最終記録までに、検証動作誤差は約0.1350から約0.1259へ低下し、相対低下は約6.7\%であった。価値推定誤差は3.1前後で変動し、単調改善ではない。\ev{E16}

\plain{強化学習は「固定した練習問題では改善した」段階であり、同一条件の実機試験はまだ行っていない。基礎モデルと強化学習モデルを、同じボトル、初期位置、回数で比較して初めて実機性能の改善幅を判断できる。}

\section{本段階で確認できる最良結果}

\begin{itemize}
  \item 把持、開栓、本体・キャップ分別投入の全工程を現場で観測した。
  \item 連続1時間超、約2本/分は現場観測による概算である。
  \item 本番運用モデルの学習、保存、運用経路を追跡できる。
  \item 学習誤差は大幅に低下したが、検証・試験・標準実機指標はない。
  \item 強化学習のオフライン動作誤差は約6.7\%改善したが、実機性能の改善幅は未確認である。
\end{itemize}

\resultbox{現段階で最も確実な成果表現は、データ生産から現場運用までの閉ループを構築し、双腕による多段階ボトル分別の全サイクルと初期的な適応兆候を現場で示した、というものである。安定成功率と標準処理量は次年度の標準試験で定量化する。}
